%%
%% fall-lecture21.tex
%% 
%% Made by Alex Nelson <pqnelson@gmail.com>
%% Login   <alex@lisp>
%% 
%% Started on  2026-11-14T08:08:47-0800
%% Last update 2026-11-14T08:08:47-0800
%% 

\lecture{}

% I'm skipping over a lot of the review
\begin{recall}
We have derived the following equations in the Shock fixed frame.
Our assumptions:
\begin{enumerate}
\item Steady flow
\item Friction is negligible
\item Heat transfer is negligible.
\end{enumerate}
First, the familiar equations of motion in integral form using a
control volume surrounding the shock wave; the continuity equation,
\begin{equation}\label{eq:me101a:fall2026:lec21:eq1}
\rho_{2}u_{2}=\rho_{1}u_{1}=\dot{m}.
\end{equation}
The conservation of momentum,
\begin{equation}\label{eq:me101a:fall2026:lec21:eq2}
p_{2}-p_{1}=\rho_{1}u_{1}^{2}-\rho_{2}u_{2}^{2}.
\end{equation}
The conservation of energy takes the form,
\begin{equation}\label{eq:me101a:fall2026:lec21:eq3}
h_{2}-h_{1}=\frac{1}{2}u_{1}^{2}-\frac{1}{2}u_{2}^{2}.
\end{equation}
Then we had
\begin{equation}\label{eq:me101a:fall2026:lec21:eq4}
\Delta p = -\dot{m}\,\Delta u,
\end{equation}
\begin{equation}\label{eq:me101a:fall2026:lec21:eq5}
\frac{\Delta u}{u_{1}}=\frac{\Delta V}{V_{1}}=\frac{-\Delta\rho}{\rho_{2}}
\end{equation}
\begin{equation}\label{eq:me101a:fall2026:lec21:eq6}
\frac{\Delta p}{\Delta\rho}=u_{1}u_{2}
\end{equation}
\begin{equation}\label{eq:me101a:fall2026:lec21:eq7}
\frac{\Delta p}{\Delta V}=-\dot{m}^{2}
\end{equation}
\begin{equation}\label{eq:me101a:fall2026:lec21:eq8}
\Delta p\,\Delta V=-(\Delta u)^{2}
\end{equation}
\begin{equation}\label{eq:me101a:fall2026:lec21:eq9}
\Delta h = \overline{V}\,\Delta p
\end{equation}
where $\overline{V}=(V_{1}+V_{2})/2$,
\begin{equation}\label{eq:me101a:fall2026:lec21:eq10}
\Delta e=-\overline{p}\,\Delta V.
\end{equation}
\end{recall}

\begin{node}[Entropy through a shock]
In quasi 1-dimensional flow, we assumed isentropic flow. This is no
longer valid for Shock waves! The second Law of Thermodyanmics takes
the form,
\begin{equation}
\underbrace{\partial_{t}\int\rho S\,\D V}_{=0}
+ \oint\rho S\vec{u}\cdot\UnitNormalVec{n}\,\D A
+ \underbrace{\oint\frac{q}{T}\D A}_{=0}
- \int_{V_{C}}\rho\frac{Q}{T}\,\D V\geq0,
\end{equation}
or (rearranging terms)
\begin{equation}
\oint\rho S\vec{u}\cdot\UnitNormalVec{n}\,\D A\geq\int_{V_{C}}\rho\frac{Q}{T}\,\D V>0
\end{equation}
since $Q\neq0$ for dissipation due to the shock (which is why the $>0$
inequality holds). Hence
\begin{equation}
\rho_{2}u_{2}S_{2}-\rho_{1}S_{1}u_{1}>0,
\end{equation}
which implies
\begin{equation}
S_{2}-S_{1}>0.
\end{equation}
So entropy always increases due to shock waves.
\end{node}

\subsection{Rayleigh Lines}

\begin{node}
From Eq~\eqref{eq:me101a:fall2026:lec21:eq7}, expanding out the differences:
\begin{equation}
\frac{p_{2}-p_{1}}{V_{2}-V_{1}}=-\dot{m}^{2},
\end{equation}
then we find
\begin{equation}
p_{2}-p_{1}=-\dot{m}^{2}(V_{2}-V_{1})
\end{equation}
where $\dot{m}$ is constant. Observe, this describes a line when we
treat $p_{2}$ and $V_{2}$ as variables. This is precisely the
\define{Rayleigh Line}.
\end{node}

\begin{node}
From Eq~\eqref{eq:me101a:fall2026:lec21:eq9},
we can write enthalpy as a function of pressure and volume $h=h(p,V)$.
This gives us a curve in the $p$-$V$ plane known as the
\define{Hugoniot Condition}.
\end{node}

\begin{node}
For a perfect gas, the Hugoniot condition may be determined more
explicitly, since
\begin{subequations}
  \begin{align}
h &= c_{p}T \\
&=\left(\frac{\gamma}{\gamma-1}\right)RT\\
&=\left(\frac{\gamma}{\gamma-1}\right)pV,
  \end{align}
\end{subequations}
then plugging this into Eq~\eqref{eq:me101a:fall2026:lec21:eq9} yields,
\begin{equation}\label{eq:me101a:fall2026:lec21:hugoniot-for-perfect-gas-eq0}
  \overline{V}\,\Delta p=\left(\frac{\gamma}{\gamma-1}\right)(p_{2}V_{2}-p_{1}V_{1}).
\end{equation}
We claim this can be rearranged to give us,
\begin{equation}\label{eq:me101a:fall2026:lec21:hugoniot-for-perfect-gas}
\boxed{\frac{p_{2}}{p_{1}}=\frac{\displaystyle{\frac{\gamma+1}{\gamma-1}-\frac{V_{2}}{V_{1}}}}{\displaystyle{\left(\frac{\gamma+1}{\gamma-1}\right)\left(\frac{V_{2}}{V_{1}}\right)-1}}.}
\end{equation}
\end{node}

\begin{proof}[Calculation ({Eq~\eqref{eq:me101a:fall2026:lec21:hugoniot-for-perfect-gas}})]
We begin by expanding out $\Delta p =p_{2}-p_{1}$ and
$\overline{V}=(V_{1}+V_{2})/2$ in Eq~\eqref{eq:me101a:fall2026:lec21:hugoniot-for-perfect-gas-eq0},
\begin{subequations}
\begin{equation}
\frac{V_{1}+V_{2}}{2}(p_{2}-p_{1})=\left(\frac{\gamma}{\gamma-1}\right)(p_{2}V_{2}-p_{1}V_{1}).
\end{equation}
Then we divide both sides by $p_{1}$,
\begin{equation}
  \frac{V_{1}+V_{2}}{2}\left(\frac{p_{2}}{p_{1}}-1\right)
  =\left(\frac{\gamma}{\gamma-1}\right)\left(\frac{p_{2}}{p_{1}}V_{2}-V_{1}\right).
\end{equation}
Then we divide both sides by $V_{1}$,
\begin{equation}
\frac{1}{2}\left(1+\frac{V_{2}}{V_{1}}\right)\left(\frac{p_{2}}{p_{1}}-1\right)=\left(\frac{\gamma}{\gamma-1}\right)\left(\frac{p_{2}}{p_{1}}\frac{V_{2}}{V_{1}}-1\right).
\end{equation}
Now, we want to move $p_{2}/p_{1}$ to one side of the equation:
\begin{equation}
  \frac{1}{2}\left(1+\frac{V_{2}}{V_{1}}\right)\left(\frac{p_{2}}{p_{1}}\right)
- \left(\frac{\gamma}{\gamma-1}\right)\frac{p_{2}}{p_{1}}\frac{V_{2}}{V_{1}}
  =-\left(\frac{\gamma}{\gamma-1}\right)
+ \frac{1}{2}\left(1+\frac{V_{2}}{V_{1}}\right)
\end{equation}
We collect terms and multiply both sides by 2
\begin{equation}
  \left[\left(1+\frac{V_{2}}{V_{1}}\right)
- 2\left(\frac{\gamma}{\gamma-1}\right)\frac{V_{2}}{V_{1}}\right]\frac{p_{2}}{p_{1}}
  =-2\left(\frac{\gamma}{\gamma-1}\right)
+ \left(1+\frac{V_{2}}{V_{1}}\right)
\end{equation}
Then using the fact that 
\begin{equation*}
1-\frac{2\gamma}{\gamma-1}=\frac{\gamma-1-2\gamma}{\gamma-1}=\frac{-1-\gamma}{-1+\gamma}=-\frac{\gamma+1}{\gamma-1},
\end{equation*}
we see that we obtain
\begin{equation}
  \left[1
- \left(\frac{\gamma+1}{\gamma-1}\right)\frac{V_{2}}{V_{1}}\right]\frac{p_{2}}{p_{1}}
  =-\left(\frac{\gamma+1}{\gamma-1}\right)
+ \left(\frac{V_{2}}{V_{1}}\right)
\end{equation}
\end{subequations}
Multiply both sides by $-1$ and isolating $p_{2}/p_{1}$ gives us the result.
\end{proof}

\begin{remark}
Conversely, we can express the Hugoniot condition for a perfect gas as:
\begin{equation}
\frac{V_{2}}{V_{1}}=\frac{\displaystyle{\frac{\gamma+1}{\gamma-1}+\frac{p_{2}}{p_{1}}}}{\displaystyle{1+\frac{\gamma+1}{\gamma-1}\frac{p_{2}}{p_{1}}}}.
\end{equation}
This is just by simple algebra. Let $x=V_{2}/V_{1}$, $y=p_{2}/p_{1}$,
and $c=(\gamma+1)/(\gamma-1)$,
then the relations are just
\begin{subequations}
  \begin{align}
    y &= \frac{c-x}{cx-1}
    \intertext{and}
    x &= \frac{c+y}{1+cy}.
  \end{align}
\end{subequations}
\end{remark}

\begin{remark}
Also remember that the classical equipartition Theorem relates
$\gamma$ to the number of thermally accessible degrees of freedom $f$
of a molecule by
\begin{subequations}
\begin{equation}
\gamma=1+\frac{2}{f},
\end{equation}
so
\begin{equation}
\frac{\gamma+1}{\gamma-1}=1+f.
\end{equation}
\end{subequations}
Monatomic gases has $f=3$, diatomic gases have $f=5$.
\end{remark}

\begin{node}
We can consider an example plot in the $p$-$V$ plane with two
isentropes describing the fluid before the shock wave $S=S_{1}$ and after the
shock wave $S=S_{2}$. Suppose $\gamma=3/2$ for simplicity. % So c=5
Then we draw the two isentropes and plot the Hugoniot condition
(dashed in red):
\begin{center}
\includegraphics{img/img.47}
\end{center}
The Hugoniot condition gives us the point on the $S_{2}$ isentrope,
and using the initial state and the intersection state we can draw a
straight line (drawn in red):
\begin{center}
\includegraphics{img/img.48}
\end{center}
This red line is the Rayleigh line. Observe that the Rayleigh line is
tangent to the $S_{2}$ isentrope.\footnote{\dots which required
plotting convex functions $p\sim 1/V$ and I hope you appreciate the
hard work it took to draw them!}
\end{node}

\begin{node}
Note that the shaded trapezoidal region in the figure:
\begin{center}
\includegraphics{img/img.49}
\end{center}
Its area is
\begin{equation}
A = \frac{p_{1}+p_{2}}{2}(V_{1}-V_{2})=-\overline{p}\,\Delta V=\Delta e
\end{equation}
which is just the change in internal energy. Similarly, if we consider
instead the trapezoid:
\begin{center}
\includegraphics{img/img.50}
\end{center}
Its area is
\begin{equation}
A = \frac{V_{1}+V_{2}}{2}(p_{2}-p_{1})=\overline{V}\,\Delta p=\Delta h.
\end{equation}
These just give us some geometric interpretation of our results.
\end{node}

\begin{node}
The Hugoniot condition for a perfect gas, Eq~\eqref{eq:me101a:fall2026:lec21:hugoniot-for-perfect-gas},
recall is:
\begin{equation*}\tag{\ref{eq:me101a:fall2026:lec21:hugoniot-for-perfect-gas}}
\frac{p_{2}}{p_{1}}=\frac{\displaystyle{\frac{\gamma+1}{\gamma-1}-\frac{V_{2}}{V_{1}}}}{\displaystyle{\left(\frac{\gamma+1}{\gamma-1}\right)\left(\frac{V_{2}}{V_{1}}\right)-1}}.
\end{equation*}
This is a curve of the form
\begin{equation}
y = \frac{c-x}{cx-1}
\end{equation}
where $x=V_{2}/V_{1}$, $y=p_{2}/p_{1}$, and $c=(\gamma+1)/(\gamma-1)$.
We can then plot such a curve:
\begin{center}
  \includegraphics{img/img.51}
\end{center}
There is an asymptote located at $V_{2}/V_{1}=(\gamma-1)/(\gamma+1)$
[i.e., $x=1/c$] and this describes the minimum volume possible
(reachable through a normal shock). Similarly, we see in the
$(p_{2}/p_{1})\to\infty$ limit,
\begin{subequations}
  \begin{align}
\frac{\rho_{2}}{\rho_{1}} &\to\frac{\gamma+1}{\gamma-1}\\
\frac{u_{2}}{u_{1}} &\to\left(\frac{u_{2}}{u_{1}}\right)_{\text{min}}=\frac{\gamma-1}{\gamma+1},\\
\frac{V_{2}}{V_{1}} &\to\left(\frac{V_{2}}{V_{1}}\right)_{\text{min}}=\frac{\gamma-1}{\gamma+1}.
  \end{align}
\end{subequations}
These describe the (upper bound on the) largest possible change due to
a normal shock. 
\end{node}

\subsection{Upstream Mach number}

\begin{node}
In the shock fixed frame, we are essentially centered on the shockwave
[drawn in thick red] and the fluid is moving at us:
\begin{center}
\includegraphics{img/img.52}
\end{center}
In these coordinates, the fluid ``upstream'' (to the left) has yet to
be ``shocked''.
In the lab frame, the shock wave is moving and the fluid upstream is stationary:
\begin{center}
\includegraphics{img/img.53}
\end{center}
The upstream Mach number is
\begin{equation}
\Mach_{1} = \frac{u_{1}}{a_{1}}.
\end{equation}
\end{node}

\begin{claim}
We have
\begin{equation}\label{eq:me101a:lec21:fall2026:claim1}
\frac{p_{2}}{p_{1}}=1+\gamma\Mach_{1}^{2}\left(1-\frac{\rho_{1}}{\rho_{2}}\right).
\end{equation}
\end{claim}

\begin{proof}
From the conservation of momentum,
\begin{equation}
\Delta p = \rho_{1}u_{1}(u_{1}-u_{2}),
\end{equation}
we obtain
\begin{subequations}
  \begin{align}
\frac{p_{2}}{p_{1}}-1 &= \frac{\rho_{1}u_{1}^{2}}{p_{1}}\left(1-\frac{u_{2}}{u_{1}}\right)\quad\mbox{by dividing through by $p_{1}$}\\
  &=\frac{\gamma u_{1}^{2}}{a_{1}^{2}}\left(1-\frac{u_{2}}{u_{1}}\right)\quad\mbox{by Eq~\eqref{eq:me101a:lec17:fall2026:speed-of-sound}}\\
  \end{align}
so
\begin{equation}
\frac{p_{2}}{p_{1}} = 1 + \frac{\gamma u_{1}^{2}}{a_{1}^{2}}\left(1-\frac{u_{2}}{u_{1}}\right).
\end{equation}
Then by mass conservation
\begin{equation}
\frac{p_{2}}{p_{1}} = 1 + \frac{\gamma u_{1}^{2}}{a_{1}^{2}}\left(1-\frac{\rho_{1}}{\rho_{2}}\right),
\end{equation}
\end{subequations}
which proves the claim.
\end{proof}

\begin{claim}
We have
\begin{equation}\label{eq:me101a:lec21:fall2026:claim2}
\frac{T_{2}}{T_{1}} = 1 + \frac{\gamma-1}{2}\Mach_{1}^{2}\left(1-\left(\frac{\rho_{1}}{\rho_{2}}\right)^{2}\right).
\end{equation}
\end{claim}

\begin{proof}
Using $h=c_{p}T$ and the conservation of enthalpy,
\begin{subequations}
\begin{equation}
c_{p}T_{1} + \frac{1}{2}u_{1}^{2}
= c_{p}T_{2} + \frac{1}{2}u_{2}^{2}
\end{equation}
so
\begin{equation}
c_{p}T_{2}
= c_{p}T_{1} + \frac{1}{2}u_{1}^{2} - \frac{1}{2}u_{2}^{2}
\end{equation}
then dividing both sides by $c_{p}T_{1}$ gives us
\begin{align}
\frac{T_{2}}{T_{1}} &= 1 + \frac{u_{1}^{2}}{2c_{p}T_{1}}\left(1-\frac{u_{2}^{2}}{u_{1}^{2}}\right)\\
&= 1 + \frac{\gamma-1}{2}\Mach_{1}^{2}\left(1-\frac{u_{2}^{2}}{u_{1}^{2}}\right)
\end{align}
since
\begin{equation}
\frac{u_{1}^{2}}{2c_{p}T_{1}} = \frac{u_{1}^{2}}{2}\left(\frac{\gamma}{\gamma-1}\right)RT_{1}=\frac{\gamma-1}{2}\Mach_{1}^{2}.
\end{equation}
\end{subequations}
Then using the conservation of mass proves the claim.
\end{proof}

\begin{node}
Now, we use the ideal gas Law
\begin{equation}
\frac{p}{\rho}=RT
\end{equation}
to consider
\begin{equation}
\frac{p_{2}/\rho_{2}}{p_{1}/\rho_{1}}=\frac{RT_{2}}{RT_{1}}=\frac{T_{2}}{T_{1}}.
\end{equation}
Then we have
\begin{subequations}
  \begin{align}
1 + \frac{\gamma-1}{2}\Mach_{1}^{2}\left(1-\left(\frac{\rho_{1}}{\rho_{2}}\right)^{2}\right)
&=\frac{T_{2}}{T_{1}}
\quad\mbox{by Eq~\eqref{eq:me101a:lec21:fall2026:claim2}}\\
&=\left(\frac{p_{2}}{p_{1}}\right)\left(\frac{\rho_{1}}{\rho_{2}}\right)\\
&=\left[1+\gamma\Mach_{1}^{2}\left(1-\frac{\rho_{1}}{\rho_{2}}\right)\right]\left(\frac{\rho_{1}}{\rho_{2}}\right)
  \end{align}
\end{subequations}
by Eq~\eqref{eq:me101a:lec21:fall2026:claim1}. Setting $x=\rho_{1}/\rho_{2}$,
we see this is precisely the quadratic equation
\begin{equation}
1 + \frac{\gamma-1}{2}\Mach_{1}^{2}(1-x^{2})=(1+\gamma\Mach_{1}^{2}(1-x))x.
\end{equation}
And we know how to solve quadratic equations!

There is the trivial solution $x=1$ which reduces to $1=1$.

Then we have the nontrivial solution
\begin{equation}
\frac{\rho_{1}}{\rho_{2}}=\frac{2+(\gamma-1)\Mach_{1}^{2}}{(\gamma+1)\Mach_{1}^{2}}.
\end{equation}
\end{node}

\begin{node}
We also see that
\begin{equation}
\frac{u_{2}}{u_{1}}=\frac{\rho_{1}}{\rho_{2}}=\frac{2+(\gamma-1)\Mach_{1}^{2}}{(\gamma+1)\Mach_{1}^{2}}.
\end{equation}
From
\begin{equation}
\frac{\Delta u}{u_{1}}=-\frac{\Delta\rho}{\rho_{2}}=\frac{\rho_{1}}{\rho_{2}}-1,
\end{equation}
we see that
\begin{subequations}
  \begin{align}
\frac{\Delta u}{a_{1}} &= \frac{u_{1}}{a_{1}}\frac{\Delta u}{u_{1}}\\
&=\Mach_{1}\frac{\Delta u}{u_{1}}\\
&=\frac{-2}{\gamma+1}\frac{\Mach_{1}^{2}-1}{\Mach_{1}}.
  \end{align}
\end{subequations}
\end{node}

\begin{node}
So this is the name of the game: we want to express quantities in
terms of the upstream Mach number $\Mach_{1}$, and that will give us a
great deal of information.
\end{node}